\documentclass[sigconf]{acmart}

\setcopyright{none}
\settopmatter{printacmref=false}
\renewcommand\footnotetextcopyrightpermission[1]{}

\acmDOI{}
\acmISBN{}
\acmConference[Conference'26]{Conference Name}{Month DD--DD, 2026}{City, Country}
\acmYear{2026}
\copyrightyear{2026}

\title{Impact of AMC-Oriented CW Adversarial Attacks on CRC Integrity}

\author{Anonymous Author}
\affiliation{
  \institution{Affiliation Withheld for Review}
  \country{N/A}
}
\email{anonymous@anonymous.org}

\begin{document}

\begin{abstract}
This paper studies a practical two-stage receiver pipeline in which automatic modulation
classification (AMC) controls downstream demodulation and CRC verification. Using the existing
\texttt{crc\_experiment.py} implementation, we run a two-track evaluation: Track A uses real
RML2016.10a IQ data to measure AWN behavior under clean and Carlini-Wagner (CW) attack
conditions; Track B uses synthetic bursts with known payload bits to measure CRC pass rates.
Perturbations are generated on Track A and transferred to Track B. Across digital modulations,
average AMC accuracy drops from 97.19\% to 4.65\% under CW. However, CRC under oracle
demodulation changes only from 63.62\% to 59.59\%. CRC collapse appears when demodulation is
driven by attacked AMC output: Adv+AMC CRC falls to 4.12\%. At 18 dB, the gap is stronger:
93.25\% (Adv+Oracle) vs. 7.81\% (Adv+AMC). The results indicate that the dominant failure path
is AMC misclassification leading to wrong demodulator selection, not only direct waveform
destruction.
\end{abstract}

\keywords{adversarial machine learning, modulation classification, RF security, CRC, CW attack}

\maketitle

\section{Introduction}
Practical receivers often use a two-stage pipeline: automatic modulation classification (AMC)
followed by modulation-specific demodulation and CRC checking. This architecture is effective
under noise, but it introduces a strong dependency: if AMC is manipulated, the receiver may
select the wrong demodulator, and CRC can fail even when the waveform is still demodulatable
under the correct modulation assumption.

This paper analyzes that failure path using the implementation in \texttt{crc\_experiment.py}
from the current codebase. We focus on three questions:
\begin{enumerate}
\item Does CW mainly break AMC, waveform demodulability, or both?
\item How large is the gap between oracle demodulation and AMC-driven demodulation under attack?
\item Are the conclusions stable across low and high SNR regimes (0 and 18 dB)?
\end{enumerate}

\section{Method}
\subsection{Two-Track Design}
The experiment uses a two-track setup.
\begin{enumerate}
\item \textbf{Track A (AMC):} real \texttt{RML2016.10a} IQ samples are classified by AWN on clean
and CW-attacked inputs to obtain AMC accuracy and perturbations.
\item \textbf{Track B (CRC):} synthetic bursts with known payload bits are generated, demodulated,
and validated with CRC pass/fail outcomes.
\end{enumerate}

Code references:
\texttt{crc\_experiment.py:73}, \texttt{crc\_experiment.py:175}, \texttt{crc\_experiment.py:222}.

\subsection{Transfer Perturbation}
CW perturbations are generated on Track A and then transferred to Track B:
\begin{equation}
x_{\text{adv,synth}} = x_{\text{synth}} + \delta_{\text{real}},
\end{equation}
where $\delta_{\text{real}} = x_{\text{adv,real}} - x_{\text{real}}$.

Code references:
\texttt{crc\_experiment.py:212}, \texttt{crc\_experiment.py:238}.

\subsection{Evaluation Scenarios}
For each modulation-SNR cell, four scenarios are evaluated:
\begin{enumerate}
\item \textbf{Clean+Oracle:} clean IQ, true modulation used for demodulation.
\item \textbf{Clean+AMC:} clean IQ, demodulation chosen from AMC behavior.
\item \textbf{Adv+Oracle:} adversarial IQ, true modulation used for demodulation.
\item \textbf{Adv+AMC:} adversarial IQ, demodulation chosen from attacked AMC output.
\end{enumerate}

The adversarial window is spliced back into full burst context before demodulation to preserve
guard-related filtering behavior (\texttt{crc\_experiment.py:93--124}).

\section{Experimental Setup}
\subsection{Model and Data}
\begin{enumerate}
\item AMC model: AWN loaded from \texttt{2016.10a\_AWN.pkl}.
\item Real data (Track A): \texttt{./data/RML2016.10a\_dict.pkl}.
\item Synthetic data (Track B): generated by \texttt{util.synth\_txrx.generate\_burst}.
\item CRC-capable digital modulations: BPSK, QPSK, 8PSK, QAM16, QAM64, PAM4, CPFSK, GFSK.
\end{enumerate}

\subsection{Attack and Output Scope}
\begin{enumerate}
\item Attack: CW via \texttt{create\_attack('cw', ...)}.
\item Track A uses 500 real samples per modulation-SNR cell.
\item Result file for this paper: \texttt{crc\_experiment\_results/crc\_vs\_amc.csv}.
\item Track B count in this result: $n=200$ synthetic bursts per digital cell.
\item SNR settings in the result file: 0 and 18 dB.
\end{enumerate}

\section{Results}
\subsection{Overall Means on Digital Modulations}
Table~\ref{tab:overall} reports means over all digital cells (8 modulations $\times$ 2 SNR
levels $\times$ 4 scenarios).

\begin{table}[t]
\caption{Overall Means on Digital Modulations}
\label{tab:overall}
\centering
\begin{tabular}{lcc}
\toprule
Scenario & AMC Acc (\%) & CRC Pass (\%) \\
\midrule
Clean+Oracle & 97.19 & 63.62 \\
Clean+AMC    & 97.19 & 62.12 \\
Adv+Oracle   & 4.65  & 59.59 \\
Adv+AMC      & 4.65  & 4.12  \\
\bottomrule
\end{tabular}
\end{table}

CW collapses AMC (97.19\% $\rightarrow$ 4.65\%). However, with oracle demodulation, CRC drops
only moderately (63.62\% $\rightarrow$ 59.59\%). CRC collapse occurs primarily in Adv+AMC.

\subsection{SNR-Stratified Means}
\begin{table}[t]
\caption{SNR-Stratified Means on Digital Modulations}
\label{tab:snr}
\centering
\begin{tabular}{lcc}
\toprule
Scenario (SNR) & AMC Acc (\%) & CRC Pass (\%) \\
\midrule
Clean+Oracle (0 dB) & 96.28 & 27.81 \\
Clean+AMC (0 dB)    & 96.28 & 26.94 \\
Adv+Oracle (0 dB)   & 2.30  & 25.94 \\
Adv+AMC (0 dB)      & 2.30  & 0.44  \\
\midrule
Clean+Oracle (18 dB) & 98.10 & 99.44 \\
Clean+AMC (18 dB)    & 98.10 & 97.31 \\
Adv+Oracle (18 dB)   & 7.00  & 93.25 \\
Adv+AMC (18 dB)      & 7.00  & 7.81  \\
\bottomrule
\end{tabular}
\end{table}

At 18 dB, the gap between Adv+Oracle and Adv+AMC CRC is 85.44 points (93.25\% vs. 7.81\%),
highlighting the AMC-induced error path.

\subsection{Per-Modulation CRC at 18 dB}
\begin{table}[t]
\caption{CRC Pass Rate at 18 dB by Modulation}
\label{tab:mod18}
\centering
\begin{tabular}{lccc}
\toprule
Modulation & Clean+Oracle & Adv+Oracle & Adv+AMC \\
\midrule
BPSK  & 100.0 & 100.0 & 39.0 \\
QPSK  & 100.0 & 100.0 & 1.5 \\
8PSK  & 100.0 & 100.0 & 0.5 \\
QAM16 & 100.0 & 99.5  & 0.5 \\
QAM64 & 96.0  & 76.0  & 0.5 \\
PAM4  & 100.0 & 71.0  & 2.0 \\
CPFSK & 100.0 & 100.0 & 18.5 \\
GFSK  & 99.5  & 99.5  & 0.0 \\
\bottomrule
\end{tabular}
\end{table}

In 13 out of 16 digital cells, the Clean+Oracle and Adv+Oracle CRC difference is within
2 percentage points. In contrast, 14 out of 16 cells have Adv+AMC CRC at or below 5\%.

\section{Discussion}
\subsection{Primary Failure Mechanism}
The dominant path is not always direct waveform destruction. Instead, the main mechanism is:
CW causes AMC misclassification, then the receiver selects the wrong demodulator, leading to
CRC failure.

\subsection{Why 18 dB Is More Diagnostic}
At 0 dB, several modulations already have low clean CRC baselines, which mixes channel-limited
errors with attack-induced errors. At 18 dB, clean baselines are near saturation, so attack
effects are easier to isolate.

\subsection{Modulation-Specific Sensitivity}
QAM64 and PAM4 show notable Adv+Oracle drops (96.0\% to 76.0\%, and 100.0\% to 71.0\%,
respectively), suggesting direct sensitivity beyond classification errors.

\section{Threats to Validity}
\begin{enumerate}
\item The study is transfer-based: perturbations are generated on Track A real data and applied
to Track B synthetic bursts. Conclusions are about transferability, not full real link attacks.
\item Clean+AMC and Adv+AMC in Track B are distribution-driven simulations based on Track A AMC
behavior, not per-sample AMC inference on synthetic bursts.
\item The CSV artifact does not fully encode all CW hyperparameters used at runtime unless
captured in command logs.
\item Analog modulations (WBFM, AM-DSB, AM-SSB) are reported without CRC metrics ($n=0$).
\end{enumerate}

\section{Conclusion}
Under the evaluated setup, CW nearly disables AMC, and CRC collapse is mainly realized through
the AMC-dependent demodulation path. This identifies a practical system risk: classification
robustness is not an isolated objective but a direct reliability dependency for link integrity.

\section*{Artifact and Reproducibility}
All numeric results in this paper come from
\texttt{crc\_experiment\_results/crc\_vs\_amc.csv}. A baseline reproduction command is:
\begin{verbatim}
python crc_experiment.py --ckpt_path ./checkpoint \
  --output_dir ./crc_experiment_results
\end{verbatim}
For exact reproducibility, explicitly set and record
\texttt{--n\_bursts}, \texttt{--cw\_c}, \texttt{--cw\_steps}, \texttt{--seed}, and
\texttt{--ta\_box}.


\begin{thebibliography}{00}

\bibitem{carlini2017}
N. Carlini and D. Wagner, ``Towards Evaluating the Robustness of Neural Networks,''
in \emph{Proc. IEEE Symposium on Security and Privacy}, 2017, pp. 39--57.

\bibitem{oshea2016}
T. J. O'Shea, J. Corgan, and T. C. Clancy,
``Convolutional Radio Modulation Recognition Networks,''
in \emph{Proc. International Conference on Engineering Applications of Neural Networks},
2016.

\bibitem{deepsig2016}
DeepSig, ``RadioML 2016.10A Dataset,'' 2016. [Online]. Available:
\url{https://www.deepsig.ai/datasets}

\bibitem{repo}
adversarial-rf project, ``AWN and adversarial RF experiments,''
repository artifact, 2026.

\end{thebibliography}


\end{document}
